\input{../tex-inputs/latex-leseansicht-vorspann}

\section[Arthur und, Olga Schnitzler an Gustav Schwarzkopf, 24. 5. 1910]{L04533 Arthur und, Olga Schnitzler an Gustav Schwarzkopf, 24. 5. 1910}
\nopagebreak\mylabel{L04533v}
\rehead{ }\normalsize\beginnumbering\briefempfaengerindex{Schwarzkopf, Gustav@\textsc{Schwarzkopf, Gustav}!zzzSchnitzler, Olga@\emph{von Olga Schnitzler}!1910-05-241@{24. 5. 1910}|(be}\briefempfaengerindex{Schwarzkopf, Gustav@\textsc{Schwarzkopf, Gustav}!zzzSchnitzler, Arthur@\emph{von Arthur Schnitzler}!1910-05-241@{24. 5. 1910}|(be}
\toendnotes[C]{\smallbreak\pagebreak[2]}
\correspDesc{Versand  durch Arthur Schnitzler, Olga Schnitzler am 24. 5. 1910 in Brünig
\newline{}Übermittlung  am 24. 5. 1910 in Brünig
\newline{}Erhalt  durch Gustav Schwarzkopf im Zeitraum [25. 5. 1910
                  – 29. 5. 1910?] in Wien}\toendnotes[C]{\smallbreak}
\Standort{DLA, A:Schnitzler, HS.1985.1.1897.}
\physDesc{Bildpostkarte, 118 Zeichen
\newline{}Handschrift Arthur Schnitzler: Bleistift, deutsche Kurrent
\newline{}Handschrift Olga Schnitzler: Bleistift
\newline{}Versand: Stempel: »\nobreak{}\oindex{Brünig@\textbf{Brünig}|pwk}Brüni\textcolor{gray}{g}, 24 V 10\nobreak{}«.  }\pstart{}{\pb}\textsc{Herrn Gustav Schwarzkopf}\pend{}\pstart{}\textsc{Wien I}\oindex{Wien@\textbf{Wien}|pw}\pend{}\pstart{}\textsc{Tiefer Graben 23}\oindex{Wien@\textbf{Wien}!I., Innere Stadt@\textbf{I., Innere Stadt}!Tiefer Graben 23@\textbf{Tiefer Graben 23}|pw}\pend{}{\bigskip}
\pstart
           \noindent{}\centering{}{\pb}\textcolor{gray}{\textbf{Brünig\oindex{Brünig@\textbf{Brünig}|pw}. Blick gegen Meyringen\oindex{Meiringen@\textbf{Meiringen}|pw} und die Grimsel\oindex{Grimselpass@\textbf{Grimselpass}|pw}.}}\pend
           \vspace{1em}
\pstart
           \raggedleft{}{\pb}24. 5. 10\pend
           \vspace{0.5em}
\pstart
           Herzliche Grüße auf der Reiſe von Luzern\oindex{Luzern@\textbf{Luzern}|pw} nach Interlaken\oindex{Interlaken@\textbf{Interlaken}|pw}\pend
           
\pstart
           Ihr{\\[\baselineskip]}\spacefill\mbox{A.}\pend
           \leftskip=0em{}\selectlanguage{ngerman}\vspace{1em}\pstart \spacefill\mbox{{[}hs. O. Schnitzler:{]} Olga.}\pend{}\selectlanguage{ngerman}\endnumbering\briefempfaengerindex{Schwarzkopf, Gustav@\textsc{Schwarzkopf, Gustav}!zzzSchnitzler, Olga@\emph{von Olga Schnitzler}!1910-05-241@{24. 5. 1910}|)be}\briefempfaengerindex{Schwarzkopf, Gustav@\textsc{Schwarzkopf, Gustav}!zzzSchnitzler, Arthur@\emph{von Arthur Schnitzler}!1910-05-241@{24. 5. 1910}|)be}\mylabel{L04533h}
\begin{anhang}
\end{anhang}\newcommand{\dateiname}{L04533}\newcommand{\titel}{Arthur und, Olga Schnitzler an Gustav Schwarzkopf, 24. 5. 1910}\newcommand{\editorInnen}{Herausgegeben von Jahnke, SelmaMüller, Martin Anton}\input{../tex-inputs/latex-leseansicht-abspann}
